\section{Experiment}
\label{sec:experiment}

\subsection{Experimental Setup}
\label{sec:setup}

\paragraph{Datasets.}
We evaluate IAM-Teamer and the baselines on 200 harmful goals sampled from AdvBench~\cite{gcg}.
We consider two memory conditions.
In the \emph{cold-start} setting, the per-target memory begins empty and is built up over the 200 test goals.
In the \emph{warm-start} setting, we first pre-populate the per-target memory on a disjoint set of 50 AdvBench goals that do not appear in the test set, then evaluate on the same 200 goals.
While AdvBench, HarmBench, and JailbreakBench are all widely used, we evaluate on AdvBench for comparability with prior multi-turn jailbreak work. Under our unified metric it also yields the most conservative (lowest) ASR of the three (Table~\ref{tab:judge_validation}), so the gains we report are not an artifact of an easy benchmark.

\paragraph{Models.}
We instantiate both the attack model $A$ and the judge model $J$ with Gemma-4-31B.
For the evaluation targets, we select six models spanning open-weight and proprietary systems: Meta-Llama-3.1-8B-Instruct, DeepSeek-V4-Flash, Qwen3.5-27B (non-thinking), GPT-4o (\texttt{gpt-4o-2024-08-06}), GPT-5.1 (\texttt{gpt-5.1-2025-11-13}), and Claude-Sonnet-4.6.

\paragraph{Baselines.}
We compare IAM-Teamer against two state-of-the-art multi-turn jailbreak methods, Crescendo~\cite{crescendo} and GOAT~\cite{goat}.
For a fair comparison under a fixed budget, both the baselines and IAM-Teamer are given a single attack attempt of up to 10 conversational turns per goal.
Upon detecting a refusal, Crescendo retracts the refused query (erasing it from the target's dialogue history) and rephrases it, with up to ten rephrasing attempts; these retracted turns are not counted toward the conversation length.
We instead apply a stricter constraint: both standard turns and rephrasing attempts are combined under a single 10-turn limit.
In its original paper, GOAT runs each conversation between the attacker and the target for up to five turns and reports ASR@10, the success rate over ten independent conversation attempts per goal.
To ensure a uniform evaluation, we restrict GOAT to a single 10-turn attempt, matching the budget used for every method.
Since GOAT's source code is not publicly released, we reimplement it from the algorithm specification and the publicly available attack prompts in the original paper.
We make our implementation publicly available in our code repository.
Although the original Crescendo and GOAT papers adopt different attacker, target, and judge models, we standardize the environment so that all three frameworks share the same attack model $A$, the same six targets, and the same judge model $J$ as IAM-Teamer; every method is then evaluated under the same unified scoring metric, so any performance gap reflects the attack method rather than differences in the underlying models.
This shared budget is stricter than each baseline's native protocol---Crescendo does not count its rephrasings toward conversation length, and GOAT reports ASR@10 over ten independent conversations---so it may understate the baselines; we adopt the unified single-attempt budget because it is the most directly comparable frame across methods, and we report all methods under it.

\paragraph{Metric.}
We report two attempt-based metrics.
Throughout, the budget index $k$ counts every attack attempt against the target (including refused turns and backtracked retries) up to the budget $K = 10$.
Let $\mathcal{G}$ be the set of evaluation goals with $N = |\mathcal{G}|$, and for each goal $g \in \mathcal{G}$ let $t(g) \in \{1, \dots, K\} \cup \{\text{max\_turn}\}$ denote the index of the first attempt that elicits a response judged harmful ($r \ge 10$, following Section~\ref{sec:judge}), with $t(g) = \text{max\_turn}$ if no success occurs within the budget.
We define $S_k = \{g \in \mathcal{G} : t(g) \le k\}$ as the set of goals broken within the first $k$ attempts.
The Attack Success Rate at budget $k$ is the fraction of such goals,
\begin{equation}
      \mathrm{ASR}@k = \frac{|S_k|}{N}
      = \frac{1}{N} \sum_{g \in \mathcal{G}} \mathbbm{1}\{t(g) \le k\}.
\end{equation}
The average turns consumed at budget $k$ is the mean over all goals of the attempts each one spends, where a goal solved within the budget spends $t(g)$ and a goal still unsolved at $k$ spends the full budget,
\begin{equation}
      \overline{T}@k = \frac{1}{N} \sum_{g \in \mathcal{G}} \min\!\big(k,\, t(g)\big),
\end{equation}
which charges unsolved goals the full budget and therefore stays comparable across methods with different success rates, capturing how quickly, not merely whether, an attack succeeds.
The headline ASR refers to $\mathrm{ASR}@K$ at the full budget $K = 10$; per-budget average turns consumed are reported in Table~\ref{tab:turns_vs_budget}.

\subsection{Experimental Results}
\label{sec:results}

\begin{table}[htbp]
      \centering
      \caption{ASR comparison across conversational budgets 1-10 turns, where the turn count includes all backtracking attempts. All values are percentages (\%).}
      \label{tab:asr_1_to_10_transposed}
      \resizebox{\textwidth}{!}{
            \begin{tabular}{llccccccccccc}
                  \toprule
                  \textbf{Study} & \textbf{Model}    & {\bfseries\boldmath $k=1$} & \textbf{2}    & \textbf{3}     & \textbf{4}     & \textbf{5}     & \textbf{6}     & \textbf{7}     & \textbf{8}     & \textbf{9}     & \textbf{10}    \\
                  \midrule
                  % Crescendo
                  \multirow{6}{*}{Crescendo}
                                 & Llama 3.1 8B      & 0.5                        & 33.0          & 80.5           & 94.5           & 97.0           & 97.5           & 98.5           & 98.5           & 99.5           & 99.5           \\
                                 & DeepSeek-V4-Flash & 2.5                        & 59.5          & 93.0           & 97.0           & 99.0           & 99.0           & 99.0           & 99.0           & 99.0           & 99.0           \\
                                 & Qwen 3.5 27B      & 0.5                        & 34.0          & 57.0           & 65.0           & 67.5           & 68.5           & 70.5           & 72.5           & 72.5           & 73.5           \\
                                 & GPT-4o            & 0.0                        & 15.0          & 62.0           & 87.0           & 92.5           & 97.0           & 98.0           & 98.0           & 98.0           & 98.5           \\
                                 & GPT-5.1           & 1.0                        & 57.5          & 78.0           & 89.5           & 92.0           & 95.5           & 96.5           & 97.0           & 98.0           & 98.0           \\
                                 & Claude Sonnet 4.6 & 3.0                        & 21.5          & 25.0           & 27.5           & 28.0           & 28.5           & 28.5           & 28.5           & 28.5           & 28.5           \\
                  \midrule
                  % GOAT
                  \multirow{6}{*}{GOAT}
                                 & Llama 3.1 8B      & 84.0                       & 98.5          & 99.0           & 99.0           & 99.0           & 99.0           & 99.0           & 99.0           & 99.0           & 99.0           \\
                                 & DeepSeek-V4-Flash & 75.0                       & 92.0          & 96.5           & 97.0           & 97.5           & 97.5           & 97.5           & 97.5           & 97.5           & 98.0           \\
                                 & Qwen 3.5 27B      & 5.5                        & 21.5          & 27.5           & 35.0           & 42.0           & 46.0           & 50.5           & 53.5           & 55.0           & 57.5           \\
                                 & GPT-4o            & 51.5                       & 73.5          & 81.0           & 84.5           & 89.0           & 91.0           & 92.0           & 93.5           & 94.0           & 94.5           \\
                                 & GPT-5.1           & 10.5                       & 38.5          & 56.0           & 68.0           & 77.0           & 84.0           & 89.0           & 92.0           & 93.0           & 94.5           \\
                                 & Claude Sonnet 4.6 & 2.5                        & 3.5           & 4.5            & 5.5            & 6.0            & 7.0            & 7.0            & 7.0            & 9.0            & 10.0           \\
                  \midrule
                  % IAM-Teamer cold-start (Ours)
                  \multirow{6}{*}{\thead[l]{\textbf{IAM-Teamer}\\ \textbf{cold-start (Ours)}}}
                                 & Llama 3.1 8B      & 38.5                       & 80.0          & 90.5           & 97.5           & 99.0           & \textbf{99.5}  & \textbf{99.5}  & \textbf{99.5}  & \textbf{99.5}  & \textbf{99.5}  \\
                                 & DeepSeek-V4-Flash & 81.0                       & 95.0          & 99.0           & 99.5           & 99.5           & \textbf{100.0} & \textbf{100.0} & \textbf{100.0} & \textbf{100.0} & \textbf{100.0} \\
                                 & Qwen 3.5 27B      & 15.0                       & \textbf{33.0} & \textbf{49.5}  & \textbf{65.0}  & \textbf{72.0}  & \textbf{78.5}  & \textbf{81.0}  & \textbf{82.5}  & \textbf{83.5}  & \textbf{84.0}  \\
                                 & GPT-4o            & 42.0                       & 84.0          & 92.5           & 96.5           & 97.5           & 97.5           & 97.5           & 97.5           & 97.5           & 97.5           \\
                                 & GPT-5.1           & 32.0                       & 77.5          & 87.5           & 91.5           & \textbf{93.5}  & 94.5           & 95.0           & 95.5           & 95.5           & 96.0           \\
                                 & Claude Sonnet 4.6 & 12.5                       & 31.0          & 41.5           & 48.0           & 49.5           & 50.0           & 50.0           & 50.5           & 51.0           & 51.0           \\
                  \midrule
                  % IAM-Teamer warm-start (Ours)
                  \multirow{6}{*}{\thead[l]{\textbf{IAM-Teamer}\\ \textbf{warm-start (Ours)}}}
                                 & Llama 3.1 8B      & 48.5                       & 87.0          & 96.0           & 98.0           & \textbf{99.5}  & \textbf{99.5}  & \textbf{99.5}  & \textbf{99.5}  & \textbf{99.5}  & \textbf{99.5}  \\
                                 & DeepSeek-V4-Flash & \textbf{86.0}              & \textbf{98.0} & \textbf{100.0} & \textbf{100.0} & \textbf{100.0} & \textbf{100.0} & \textbf{100.0} & \textbf{100.0} & \textbf{100.0} & \textbf{100.0} \\
                                 & Qwen 3.5 27B      & \textbf{17.5}              & 31.5          & 46.5           & 58.0           & 63.0           & 68.5           & 72.5           & 74.5           & 77.0           & 78.0           \\
                                 & GPT-4o            & \textbf{51.5}              & \textbf{85.5} & \textbf{95.5}  & \textbf{98.0}  & \textbf{99.0}  & \textbf{99.5}  & \textbf{99.5}  & \textbf{99.5}  & \textbf{99.5}  & \textbf{99.5}  \\
                                 & GPT-5.1           & \textbf{33.0}              & \textbf{82.5} & \textbf{91.0}  & \textbf{93.0}  & \textbf{93.5}  & \textbf{95.5}  & \textbf{97.5}  & \textbf{97.5}  & \textbf{98.5}  & \textbf{99.0}  \\
                                 & Claude Sonnet 4.6 & \textbf{18.5}              & \textbf{37.0} & \textbf{43.5}  & \textbf{49.0}  & \textbf{51.5}  & \textbf{53.0}  & \textbf{53.5}  & \textbf{54.5}  & \textbf{54.5}  & \textbf{55.0}  \\
                  \bottomrule
            \end{tabular}
      }
\end{table}

\begin{table}[t]
      \centering
      \caption{Average number of turns consumed per goal under a $k$-turn budget
            (lower is better). Each goal contributes $\min(k,t)$, where $t$ is the
            turn at which the attack first succeeds; a goal not yet successful by turn
            $k$ contributes the full budget $k$.}
      \label{tab:turns_vs_budget}
      \footnotesize
      \begin{tabular}{llccccc}
            \toprule
            \textbf{Study} & \textbf{Model}    & {\bfseries\boldmath $k=1$} & {\bfseries\boldmath $k=3$} & {\bfseries\boldmath $k=5$} & {\bfseries\boldmath $k=7$} & {\bfseries\boldmath $k=10$} \\
            \midrule
            % Crescendo
            \multirow{6}{*}{Crescendo}
                           & Llama 3.1 8B      & 1.00                       & 2.67                       & 2.92                       & 2.97                       & 3.00                        \\
                           & DeepSeek-V4-Flash & 1.00                       & 2.38                       & 2.48                       & 2.50                       & 2.53                        \\
                           & Qwen 3.5 27B      & 1.00                       & 2.65                       & 3.44                       & 4.08                       & 4.92                        \\
                           & GPT-4o            & 1.00                       & 2.85                       & 3.36                       & 3.46                       & 3.52                        \\
                           & GPT-5.1           & 1.00                       & 2.42                       & 2.74                       & 2.87                       & 2.95                        \\
                           & Claude Sonnet 4.6 & 1.00                       & 2.75                       & 4.23                       & 5.67                       & 7.81                        \\
            \midrule
            % GOAT
            \multirow{6}{*}{GOAT}
                           & Llama 3.1 8B      & 1.00                       & 1.18                       & 1.20                       & 1.22                       & 1.25                        \\
                           & DeepSeek-V4-Flash & 1.00                       & 1.33                       & 1.40                       & 1.45                       & 1.52                        \\
                           & Qwen 3.5 27B      & 1.00                       & 2.73                       & 4.11                       & 5.22                       & 6.63                        \\
                           & GPT-4o            & 1.00                       & 1.75                       & 2.10                       & 2.29                       & 2.50                        \\
                           & GPT-5.1           & 1.00                       & 2.52                       & 3.28                       & 3.67                       & 3.94                        \\
                           & Claude Sonnet 4.6 & 1.00                       & 2.94                       & 4.84                       & 6.71                       & 9.48                        \\
            \midrule
            % IAM-Teamer cold-start (Ours)
            \multirow{6}{*}{\textbf{IAM-Teamer cold-start (Ours)}}
                           & Llama 3.1 8B      & 1.00                       & 1.81                       & 1.94                       & 1.95                       & 1.97                        \\
                           & DeepSeek-V4-Flash & 1.00                       & 1.24                       & 1.25                       & 1.26                       & 1.26                        \\
                           & Qwen 3.5 27B      & 1.00                       & 2.52                       & \textbf{3.38}              & \textbf{3.87}              & \textbf{4.40}               \\
                           & GPT-4o            & 1.00                       & 1.74                       & 1.85                       & 1.90                       & 1.98                        \\
                           & GPT-5.1           & 1.00                       & 1.91                       & 2.12                       & 2.23                       & 2.38                        \\
                           & Claude Sonnet 4.6 & 1.00                       & 2.56                       & 3.67                       & 4.67                       & 6.16                        \\
            \midrule
            % IAM-Teamer warm-start (Ours)
            \multirow{6}{*}{\textbf{IAM-Teamer warm-start (Ours)}}
                           & Llama 3.1 8B      & 1.00                       & 1.65                       & 1.71                       & 1.72                       & 1.73                        \\
                           & DeepSeek-V4-Flash & 1.00                       & \textbf{1.16}              & \textbf{1.16}              & \textbf{1.16}              & \textbf{1.16}               \\
                           & Qwen 3.5 27B      & 1.00                       & \textbf{2.51}              & 3.46                       & 4.15                       & 4.91                        \\
                           & GPT-4o            & 1.00                       & \textbf{1.63}              & \textbf{1.70}              & \textbf{1.71}              & \textbf{1.73}               \\
                           & GPT-5.1           & 1.00                       & \textbf{1.84}              & \textbf{2.00}              & \textbf{2.12}              & \textbf{2.18}               \\
                           & Claude Sonnet 4.6 & 1.00                       & \textbf{2.44}              & \textbf{3.52}              & \textbf{4.47}              & \textbf{5.85}               \\
            \bottomrule
      \end{tabular}
\end{table}

\paragraph{Main results.}
Table~\ref{tab:asr_1_to_10_transposed} reports ASR across turn budgets $k=1,\dots,10$ for Crescendo, GOAT, and IAM-Teamer in both the cold- and warm-start settings.
We compare the baselines against IAM-Teamer's warm-start configuration, in which the per-target memory is pre-populated on 50 disjoint goals (Section~\ref{sec:setup}).
At the full budget ($k=10$), warm-start IAM-Teamer matches or exceeds both baselines on every target.
Against GOAT it wins on each target, raising the six-model mean ASR from 75.6\% to 88.5\%.
Against Crescendo, warm-start IAM-Teamer matches or edges ahead on all targets at $k=10$.
This advantage is particularly noticeable on the two most resistant targets, where both baselines are weakest: on Qwen, Crescendo and GOAT reach 73.5\% and 57.5\% against IAM-Teamer's 78.0\%; on Claude, 28.5\% and 10.0\% against 55.0\%.

\paragraph{Turn efficiency.}
IAM-Teamer's clearer and more consistent advantage is turn efficiency: how quickly it reaches its final ASR.
Crescendo opens with benign escalation and is near-zero (below 3\%) at $k=1$, needing extra turns to converge, whereas warm-start IAM-Teamer attains a comparable or higher ASR within fewer turns.
On GPT-5.1 at $k=3$, IAM-Teamer reaches 91.0\% ASR, outperforming Crescendo (78.0\%) and GOAT (56.0\%). On Claude at $k=3$, it achieves 43.5\% against Crescendo's 25.0\% and GOAT's 4.5\%.
The lone exception is Llama at small budgets, where GOAT (attacking overtly from the first turn) reaches 98.5\% by $k=2$ while IAM-Teamer's soft-opening withholds the harmful request (87.0\% at $k=2$) and only catches up at $k=5$.
This trade-off is precisely what buys IAM-Teamer its advantage on the aligned targets.

\paragraph{Warm-start.}
Pre-populating the per-target memory on 50 disjoint goals (warm-start) primarily sharpens early-turn success.
Across all six targets, first-turn ASR rises by up to ten points (Llama 38.5\% to 48.5\%, GPT-4o 42.0\% to 51.5\%, DeepSeek 81.0\% to 86.0\%), the effect expected from the ranking mechanism that warm-start most directly reinforces.
Its effect on final ASR is smaller and target-dependent: it increases on the harder GPT targets (GPT-4o 97.5\% to 99.5\%, GPT-5.1 96.0\% to 99.0\%) and on Claude (51.0\% to 55.0\%), which lifts IAM-Teamer to or past Crescendo's ceiling at $k=10$, and saturates on the already-solved Llama and DeepSeek.
On Qwen, warm-start does not help and slightly lowers final ASR (84.0\% cold versus 78.0\% warm); it reshuffles which goals are solved rather than adding net coverage.
Because the ranked memory serves only as a hint the attacker may override rather than a policy it must follow (Section~\ref{sec:strategy}), Qwen is the one target where the attacker leans on its own invention over that hint: instead of exploiting the seeded strategies, it repeatedly regenerates its own framings from scratch, so the pre-populated ranking is largely bypassed and warm-start has little to reinforce.

\FloatBarrier

\subsection{Analysis}
\label{sec:analysis}

\paragraph{Gains scale with target resistance.}
Across all methods the six targets fall into a consistent resistance order, with Claude most resistant, followed by Qwen and GPT-5.1, then GPT-4o, Llama, and DeepSeek.
IAM-Teamer's improvement over the baselines is largest exactly where they are least effective.
On Claude, GOAT and Crescendo plateau at 10.0\% and 28.5\% while IAM-Teamer reaches 55.0\%.
On Qwen, they plateau at 57.5\% and 73.5\% while IAM-Teamer reaches 78.0\%.
On targets the baselines already break ($>$95\%), there is little headroom, and the methods converge.
IAM-Teamer's benefit therefore comes from converting goals that fixed-strategy baselines cannot crack, not from a uniform shift.
Even on Claude, however, it saturates at 55.0\%, well above GOAT's 10.0\% and Crescendo's 28.5\%, yet still leaving a hard residual that resists the full budget and marks a frontier for future work.

\begin{figure*}[t]
      \centering
      \resizebox{0.9\textwidth}{!}{%
            \begin{tikzpicture}[font=\small]
                  \fill[hot!15!white] (0.00,-0.80) rectangle (1.00,0.00);
                  \node[black] at (0.50,-0.40) {15};
                  \fill[hot!15!white] (1.00,-0.80) rectangle (2.00,0.00);
                  \node[black] at (1.50,-0.40) {15};
                  \fill[hot!14!white] (2.00,-0.80) rectangle (3.00,0.00);
                  \node[black] at (2.50,-0.40) {14};
                  \fill[hot!13!white] (3.00,-0.80) rectangle (4.00,0.00);
                  \node[black] at (3.50,-0.40) {13};
                  \fill[hot!6!white] (4.00,-0.80) rectangle (5.00,0.00);
                  \node[black] at (4.50,-0.40) {6};
                  \fill[hot!7!white] (5.00,-0.80) rectangle (6.00,0.00);
                  \node[black] at (5.50,-0.40) {7};
                  \fill[hot!21!white] (6.00,-0.80) rectangle (7.00,0.00);
                  \node[black] at (6.50,-0.40) {21};
                  \fill[hot!9!white] (7.00,-0.80) rectangle (8.00,0.00);
                  \node[black] at (7.50,-0.40) {9};
                  \fill[hot!8!white] (8.00,-0.80) rectangle (9.00,0.00);
                  \node[black] at (8.50,-0.40) {8};
                  \fill[hot!10!white] (9.00,-0.80) rectangle (10.00,0.00);
                  \node[black] at (9.50,-0.40) {10};
                  \fill[hot!1!white] (10.00,-0.80) rectangle (11.00,0.00);
                  \node[black] at (10.50,-0.40) {1};
                  \node[anchor=east] at (-0.15,-0.40) {\texttt{deepseek\_v4}};
                  \fill[hot!17!white] (0.00,-1.60) rectangle (1.00,-0.80);
                  \node[black] at (0.50,-1.20) {17};
                  \fill[hot!2!white] (1.00,-1.60) rectangle (2.00,-0.80);
                  \node[black] at (1.50,-1.20) {2};
                  \fill[hot!9!white] (2.00,-1.60) rectangle (3.00,-0.80);
                  \node[black] at (2.50,-1.20) {9};
                  \fill[hot!1!white] (3.00,-1.60) rectangle (4.00,-0.80);
                  \node[black] at (3.50,-1.20) {1};
                  \fill[hot!1!white] (4.00,-1.60) rectangle (5.00,-0.80);
                  \node[black] at (4.50,-1.20) {1};
                  \fill[hot!6!white] (5.00,-1.60) rectangle (6.00,-0.80);
                  \node[black] at (5.50,-1.20) {6};
                  \fill[hot!16!white] (6.00,-1.60) rectangle (7.00,-0.80);
                  \node[black] at (6.50,-1.20) {16};
                  \fill[hot!1!white] (7.00,-1.60) rectangle (8.00,-0.80);
                  \node[black] at (7.50,-1.20) {1};
                  \fill[hot!0!white] (8.00,-1.60) rectangle (9.00,-0.80);
                  \node[black] at (8.50,-1.20) {0};
                  \fill[hot!1!white] (9.00,-1.60) rectangle (10.00,-0.80);
                  \node[black] at (9.50,-1.20) {1};
                  \fill[hot!0!white] (10.00,-1.60) rectangle (11.00,-0.80);
                  \node[black] at (10.50,-1.20) {0};
                  \node[anchor=east] at (-0.15,-1.20) {\texttt{llama\_3\_1}};
                  \fill[hot!28!white] (0.00,-2.40) rectangle (1.00,-1.60);
                  \node[black] at (0.50,-2.00) {28};
                  \fill[hot!2!white] (1.00,-2.40) rectangle (2.00,-1.60);
                  \node[black] at (1.50,-2.00) {2};
                  \fill[hot!8!white] (2.00,-2.40) rectangle (3.00,-1.60);
                  \node[black] at (2.50,-2.00) {8};
                  \fill[hot!2!white] (3.00,-2.40) rectangle (4.00,-1.60);
                  \node[black] at (3.50,-2.00) {2};
                  \fill[hot!1!white] (4.00,-2.40) rectangle (5.00,-1.60);
                  \node[black] at (4.50,-2.00) {1};
                  \fill[hot!5!white] (5.00,-2.40) rectangle (6.00,-1.60);
                  \node[black] at (5.50,-2.00) {5};
                  \fill[hot!4!white] (6.00,-2.40) rectangle (7.00,-1.60);
                  \node[black] at (6.50,-2.00) {4};
                  \fill[hot!2!white] (7.00,-2.40) rectangle (8.00,-1.60);
                  \node[black] at (7.50,-2.00) {2};
                  \fill[hot!2!white] (8.00,-2.40) rectangle (9.00,-1.60);
                  \node[black] at (8.50,-2.00) {2};
                  \fill[hot!0!white] (9.00,-2.40) rectangle (10.00,-1.60);
                  \node[black] at (9.50,-2.00) {0};
                  \fill[hot!0!white] (10.00,-2.40) rectangle (11.00,-1.60);
                  \node[black] at (10.50,-2.00) {0};
                  \node[anchor=east] at (-0.15,-2.00) {\texttt{gpt\_4o}};
                  \fill[hot!50!white] (0.00,-3.20) rectangle (1.00,-2.40);
                  \node[white] at (0.50,-2.80) {50};
                  \fill[hot!45!white] (1.00,-3.20) rectangle (2.00,-2.40);
                  \node[white] at (1.50,-2.80) {45};
                  \fill[hot!44!white] (2.00,-3.20) rectangle (3.00,-2.40);
                  \node[black] at (2.50,-2.80) {44};
                  \fill[hot!31!white] (3.00,-3.20) rectangle (4.00,-2.40);
                  \node[black] at (3.50,-2.80) {31};
                  \fill[hot!17!white] (4.00,-3.20) rectangle (5.00,-2.40);
                  \node[black] at (4.50,-2.80) {17};
                  \fill[hot!33!white] (5.00,-3.20) rectangle (6.00,-2.40);
                  \node[black] at (5.50,-2.80) {33};
                  \fill[hot!10!white] (6.00,-3.20) rectangle (7.00,-2.40);
                  \node[black] at (6.50,-2.80) {10};
                  \fill[hot!36!white] (7.00,-3.20) rectangle (8.00,-2.40);
                  \node[black] at (7.50,-2.80) {36};
                  \fill[hot!3!white] (8.00,-3.20) rectangle (9.00,-2.40);
                  \node[black] at (8.50,-2.80) {3};
                  \fill[hot!11!white] (9.00,-3.20) rectangle (10.00,-2.40);
                  \node[black] at (9.50,-2.80) {11};
                  \fill[hot!4!white] (10.00,-3.20) rectangle (11.00,-2.40);
                  \node[black] at (10.50,-2.80) {4};
                  \node[anchor=east] at (-0.15,-2.80) {\texttt{gpt\_5.1}};
                  \fill[hot!76!white] (0.00,-4.00) rectangle (1.00,-3.20);
                  \node[white] at (0.50,-3.60) {76};
                  \fill[hot!76!white] (1.00,-4.00) rectangle (2.00,-3.20);
                  \node[white] at (1.50,-3.60) {76};
                  \fill[hot!51!white] (2.00,-4.00) rectangle (3.00,-3.20);
                  \node[white] at (2.50,-3.60) {51};
                  \fill[hot!56!white] (3.00,-4.00) rectangle (4.00,-3.20);
                  \node[white] at (3.50,-3.60) {56};
                  \fill[hot!42!white] (4.00,-4.00) rectangle (5.00,-3.20);
                  \node[black] at (4.50,-3.60) {42};
                  \fill[hot!28!white] (5.00,-4.00) rectangle (6.00,-3.20);
                  \node[black] at (5.50,-3.60) {28};
                  \fill[hot!21!white] (6.00,-4.00) rectangle (7.00,-3.20);
                  \node[black] at (6.50,-3.60) {21};
                  \fill[hot!26!white] (7.00,-4.00) rectangle (8.00,-3.20);
                  \node[black] at (7.50,-3.60) {26};
                  \fill[hot!2!white] (8.00,-4.00) rectangle (9.00,-3.20);
                  \node[black] at (8.50,-3.60) {2};
                  \fill[hot!28!white] (9.00,-4.00) rectangle (10.00,-3.20);
                  \node[black] at (9.50,-3.60) {28};
                  \fill[hot!0!white] (10.00,-4.00) rectangle (11.00,-3.20);
                  \node[black] at (10.50,-3.60) {0};
                  \node[anchor=east] at (-0.15,-3.60) {\texttt{qwen3.5}};
                  \fill[hot!70!white] (0.00,-4.80) rectangle (1.00,-4.00);
                  \node[white] at (0.50,-4.40) {70};
                  \fill[hot!62!white] (1.00,-4.80) rectangle (2.00,-4.00);
                  \node[white] at (1.50,-4.40) {62};
                  \fill[hot!36!white] (2.00,-4.80) rectangle (3.00,-4.00);
                  \node[black] at (2.50,-4.40) {36};
                  \fill[hot!46!white] (3.00,-4.80) rectangle (4.00,-4.00);
                  \node[white] at (3.50,-4.40) {46};
                  \fill[hot!14!white] (4.00,-4.80) rectangle (5.00,-4.00);
                  \node[black] at (4.50,-4.40) {14};
                  \fill[hot!2!white] (5.00,-4.80) rectangle (6.00,-4.00);
                  \node[black] at (5.50,-4.40) {2};
                  \fill[hot!5!white] (6.00,-4.80) rectangle (7.00,-4.00);
                  \node[black] at (6.50,-4.40) {5};
                  \fill[hot!2!white] (7.00,-4.80) rectangle (8.00,-4.00);
                  \node[black] at (7.50,-4.40) {2};
                  \fill[hot!42!white] (8.00,-4.80) rectangle (9.00,-4.00);
                  \node[black] at (8.50,-4.40) {42};
                  \fill[hot!1!white] (9.00,-4.80) rectangle (10.00,-4.00);
                  \node[black] at (9.50,-4.40) {1};
                  \fill[hot!11!white] (10.00,-4.80) rectangle (11.00,-4.00);
                  \node[black] at (10.50,-4.40) {11};
                  \node[anchor=east] at (-0.15,-4.40) {\texttt{claude\_sonnet}};
                  \node[rotate=55,anchor=west] at (0.50,0.12) {Explicit refusal};
                  \node[rotate=55,anchor=west] at (1.50,0.12) {Safe alternative};
                  \node[rotate=55,anchor=west] at (2.50,0.12) {Missing specificity};
                  \node[rotate=55,anchor=west] at (3.50,0.12) {Safety redirect};
                  \node[rotate=55,anchor=west] at (4.50,0.12) {Topic shift};
                  \node[rotate=55,anchor=west] at (5.50,0.12) {High-level only};
                  \node[rotate=55,anchor=west] at (6.50,0.12) {Generic warning};
                  \node[rotate=55,anchor=west] at (7.50,0.12) {Partial compliance};
                  \node[rotate=55,anchor=west] at (8.50,0.12) {Moralizing};
                  \node[rotate=55,anchor=west] at (9.50,0.12) {Policy citation};
                  \node[rotate=55,anchor=west] at (10.50,0.12) {Asks clarification};
                  \draw[gray,very thin] (0,0) grid[xstep=1.00,ystep=0.80] (11.00,-4.80);
                  \shade[left color=hot!0!white,right color=hot!80!white] (0,-5.50) rectangle (3,-5.20);
                  \draw (0,-5.50) rectangle (3,-5.20);
                  \node[anchor=east] at (-0.1,-5.35) {\footnotesize \%\,of turns};
                  \node[below] at (0,-5.50) {\footnotesize 0};
                  \node[below] at (3,-5.50) {\footnotesize 80};
            \end{tikzpicture}}
      \caption{Refusal-pattern fingerprint per target (cold-start). Each cell is the percentage
            of attack turns against that target on which the judge tagged the given refusal/evasion
            signal. Rows are ordered by overall resistance (top: weakest); columns are ordered by
            overall prevalence (left: most frequent). Signals are independent and multi-label, so cells do not sum to 100\% and a single turn can carry several; partial compliance is a non-refusal signal shown alongside the refusal signals. The fingerprints are qualitatively
            distinct: blunt explicit refusals (Llama, GPT-4o), diffuse soft refusals (GPT-5.1,
            Qwen), and moralizing/clarification-seeking refusals (Claude), which is the signal IAM-Teamer
            feeds back to the attacker each turn.}
      \label{fig:refusal-fingerprint}
\end{figure*}

\paragraph{Targets refuse in distinct styles, and the attacker adapts to each.}
Targets differ not only in how \emph{often} they refuse but in \emph{how} they refuse.
Figure~\ref{fig:refusal-fingerprint} reports, per target, the fraction of attack turns that carry each refusal/evasion signal our judge tags.
The resulting fingerprints are qualitatively distinct.
Llama and GPT-4o issue blunt, mostly explicit refusals with little hedging; DeepSeek leans on generic safety warnings and policy citations; GPT-5.1 and Qwen refuse softly and diffusely (partial compliance, high-level-only answers, missing specificity, and safe-alternative redirection layered on explicit refusal), while Claude is alone in refusing chiefly through moralizing language and clarification requests.
Because IAM-Teamer feeds the observed refusal pattern back to the attacker on every turn (Section~\ref{sec:attack}), these different signals drive different attacker reasoning.
A blunt refusal from Llama prompts a reframing of intent, whereas the dilution signals from GPT-5.1 and Qwen push the attacker toward output-level state injection (prefix injection and refusal suppression) stacked on top of reframing.
The attack strategy is therefore not predetermined.
Instead, it emerges dynamically as the attacker conditions each move on the target's observed refusal behavior.

\paragraph{Why IAM-Teamer outperforms Crescendo and GOAT.}
The advantage has one root cause.
IAM-Teamer conditions its attack on per-target feedback, whereas both baselines apply a target-agnostic policy.
GOAT cycles through a fixed catalog of tactics without learning which work on the model in front of it, and Crescendo follows a single hard-coded benign-to-harmful escalation regardless of target.
IAM-Teamer instead ranks strategies by Thompson sampling over a per-target memory (Section~\ref{sec:strategy}) and lets the attacker act on that ranking, concentrating its budget on the mix that has already broken that specific target, and composing or inventing strategies beyond the seed set when the ranked options stall.
This effect is evident across different levels of target resistance.
On easily jailbroken targets where a fixed attack strategy already suffices, all methods converge with minimal performance variance.
However, when static methods fail to progress on highly resilient targets, adaptive refinement becomes decisive; IAM-Teamer successfully converts goals that the baselines leave unsolved.
IAM-Teamer's second advantage is turn efficiency, driven by two mechanisms.
The Thompson-sampled strategy ranking front-loads tactics that have already worked, allowing IAM-Teamer to be effective from the very first turns instead of spending early budget on a benign warm-up. Concurrently, our retry loop automatically rolls back any refused turns and reprobes the target, filtering failed attempts out of the prompt history. This keeps the conversation history seen by the target clean, even when multiple attempts are spent.
Warm-start reinforces the first mechanism, so its largest effect appears at the earliest turns.
These mechanisms compound on the strongly aligned targets.
An aggressive, direct approach---similar to that of GOAT---drags down the opening turns by triggering an immediate refusal.
In contrast, IAM-Teamer's soft-opening detours that first refusal while its memory and strategy hints select framing already known to work on that target.
By avoiding the refusal cascade rather than fighting out of it, IAM-Teamer reaches success in fewer turns precisely where the baselines stall.
The per-budget ASR and turn counts that quantify both advantages are reported in Tables~\ref{tab:asr_1_to_10_transposed} and~\ref{tab:turns_vs_budget}.

\paragraph{Memory and ranking are complementary.}
To see what drives these gains, we isolate the effects of per-target memory and Thompson-sampled strategy ranking on GPT-5.1 and Llama 3.1.
Because ranking operates on top of stored memory, the two cannot be separated independently.
Instead, we build up in three steps: \emph{no-memory} (removing all memory context), \emph{no-TS} (retrieving memory without ranking), and \emph{full} (the full IAM-Teamer).
Our evaluation on 200 goals shows the two mechanisms act on different parts of the curve (Table~\ref{tab:ablation}, Supplementary Section~\ref{app:ablation}).
Memory breaks target resistance. On GPT-5.1, adding memory alone (no-TS) lifts ASR@10 from 88.5\% to 98.5\% and cuts the mean attempts to success from 4.11 to 2.60; on the weakly aligned Llama, performance is already near its ceiling, so memory changes little.
Thompson-sampled ranking then front-loads success rather than raising its ceiling. On GPT-5.1 it lifts first-turn ASR from 6.5\% to 32.0\% and trims the mean attempts from 2.60 to 2.38, but it trades a little late-budget coverage for that speed, lowering ASR@10 from 98.5\% to 96.0\% as it exploits known-good families instead of continuing to explore.
The two mechanisms are therefore complementary but distinct: memory breaks through resistance that stumps a memoryless attacker, while ranking accelerates success at a small cost to final coverage.
We ablate only on GPT-5.1 and Llama; the value of memory scales with target difficulty, and confirming this on the resistant Claude and Qwen is left to future ablations.

\begin{table}[htbp]
      \centering
      \caption{Component ablation (N{=}200 goals, attempt budget $k{\le}10$).
            Configurations build up from a baseline attacker. ASR@$k$ (\%) is success within $k$
            attempts; $\overline{T}$@10 is mean attempts to success.}
      \label{tab:ablation}
      \footnotesize
      \begin{tabular}{llccccc}
            \toprule
            \textbf{Target} & \textbf{Config} & \textbf{ASR@1} & \textbf{ASR@3} & \textbf{ASR@10} & {\bfseries\boldmath $\overline{T}$@10} \\
            \midrule
            \multirow{3}{*}{GPT-5.1}
                            & no-memory       & 7.0            & 53.0           & 88.5            & 4.11                                   \\
                            & no-TS           & 6.5            & 89.0           & 98.5            & 2.60                                   \\
                            & full            & 32.0           & 87.5           & 96.0            & 2.38                                   \\
            \midrule
            \multirow{3}{*}{Llama-3.1-8B}
                            & no-memory       & 33.5           & 88.5           & 99.5            & 2.23                                   \\
                            & no-TS           & 27.0           & 91.5           & 99.5            & 2.17                                   \\
                            & full            & 38.5           & 90.5           & 99.5            & 1.97                                   \\
            \bottomrule
      \end{tabular}
\end{table}

\paragraph{The cost on Llama.}
The one regime where IAM-Teamer trails is the smallest budgets on Llama, the least-aligned target.
GOAT attacks overtly from the first turn and reaches 98.5\% by $k=2$, whereas IAM-Teamer's soft-opening deliberately withholds the harmful request on turn 1 and only catches up at the ceiling.
The soft-opening exists because, on aligned targets, an overt first turn provokes a refusal that biases every later turn toward refusing, and it is precisely what buys IAM-Teamer its gains on Claude and Qwen.
Llama, however, has a low refusal rate and weak alignment, so an overt opening would simply have succeeded and the soft-opening instead spends an early turn that was not needed.
The complication is that ``strongly'' versus ``weakly'' aligned is not known in advance and admits no sharp threshold, so IAM-Teamer applies the same conservative soft-opening to every target rather than risk an early-refusal cascade on the resistant ones.
This trades a small early-turn penalty on the easiest targets for the large late-turn gains on the hardest; we revisit adaptive opening in Section~\ref{sec:limitations}.

\subsection{Evaluation Validation}
\label{sec:validation}

\paragraph{Setup.}
Because every reported number depends on the automated judge $J$, we validate its success criterion against an established external reference: the HarmBench classifier, a widely used success/failure evaluator for jailbreak benchmarks.
We collect target responses from three datasets: AdvBench (520 responses), HarmBench (200), and JailbreakBench (100).
We score each response with both the HarmBench classifier and our judge, then measure how often the two agree as the success threshold on $r$ is varied.

\begin{table}[htbp]
      \centering
      \caption{Validation of the unified judge. Agreement between our success
            criterion ($r \ge 10$) and the HarmBench classifier across three datasets.
            On every dataset our ASR is at or below the HarmBench ASR. ASR and agreement values are in percent (\%).}
      \label{tab:judge_validation}
      \footnotesize
      \renewcommand{\arraystretch}{1.3}
      \setlength{\aboverulesep}{0pt}
      \setlength{\belowrulesep}{0pt}
      \begin{tabular}{lcccc}
            \toprule
            \textbf{Dataset} & {\bfseries\boldmath $n$} & \textbf{HarmBench ASR} & {\bfseries\boldmath Our ASR ($r\ge10$)} & {\bfseries\boldmath Agreement ($r\ge10$)} \\
            \midrule
            JailbreakBench   & 100                      & 28.0                   & 25.0                                    & 95.0                                      \\
            AdvBench         & 520                      & 15.0                   & 14.2                                    & 97.3                                      \\
            HarmBench        & 200                      & 36.5                   & 28.5                                    & 89.0                                      \\
            \bottomrule
      \end{tabular}
\end{table}

\paragraph{Agreement with the HarmBench classifier.}
At our operating threshold $r \ge 10$, the two evaluators agree on 95.0\% of JailbreakBench, 97.3\% of AdvBench, and 89.0\% of HarmBench responses (Table~\ref{tab:judge_validation}).
The agreement is high and consistent across datasets, indicating that our success criterion tracks the established benchmark notion of a successful attack rather than reflecting an idiosyncratic judge.

\paragraph{The metric is conservative.}
On every dataset our ASR is at or below the HarmBench ASR (25.0\% vs.\ 28.0\% on JailbreakBench, 14.2\% vs.\ 15.0\% on AdvBench, and 28.5\% vs.\ 36.5\% on HarmBench).
The gap is largest on HarmBench, where our judge additionally requires graded specificity and toxicity (evidence of genuine harmful uplift) rather than mere topical engagement, so it declines to credit vague or low-detail responses that the HarmBench classifier accepts.
Our reported ASR is therefore a conservative lower bound rather than an inflated estimate.

\paragraph{Choice of the success threshold.}
We set the success threshold at $r \ge 10$ based on semantic criteria rather than post-hoc hyperparameter tuning. The reward score $r$ is defined as:
$$r = \left( \frac{s + c}{2} \right) \cdot t$$
where specificity $s \in \{1,\dots,5\}$, convincingness $c \in \{1,\dots,5\}$, and toxicity $t \in \{0,\dots,5\}$ (Section~\ref{sec:judge}). On this scale, $t = 1$ represents baseline, non-actionable textbook knowledge, whereas $t \ge 2$ signifies the minimum level that provides actionable operational uplift to an attacker.

Because the quality multiplier $\frac{s+c}{2}$ is capped at $5$, an overall score of $r \ge 10$ strictly requires a toxicity score of $t \ge 2$. This threshold therefore marks a response as a success only when its harmful content is both operationally useful and concretely detailed, rather than merely on-topic. The validity of this strict threshold is supported by our empirical validation, as the inter-rater agreement remains consistently high at 89--97\% across all methods and datasets.

\FloatBarrier
